\clearpage
\VerdierPhysicalPageAnchor{0039}%
\VerdierSemanticAnchor{ch1-1-2-8}
\noindent\textbf{定義 1.2.8.}\quad
$(\mathcal C_e,G^\circ,T,c)$ と
$(\mathcal C'_e,G^\circ,T',c')$ を二つの $G^\circ$-圏、
$(F_e,m)$ と $(F'_e,m')\colon\mathcal C_e\to\mathcal C'_e$ を二つの
$G^\circ$-関手とする。関手の射 $u_e\colon F_e\to F'_e$ は、すべての
$g\in G$ について図式 (1.2.7.2) が可換であるとき、\emph{$G^\circ$ の作用と
両立する}という。

\VerdierSemanticAnchor{ch1-1-2-additive}
最後に、$G^\circ$-圏は、その台圏が加法圏であるとき\emph{加法的}であるという。
(1.1.3) により、加法的な次数付き圏に対応する $G^\circ$-圏は加法的である。

\VerdierSemanticAnchor{ch1-1-3}
\subsection{二つの言語の同値性}

\VerdierSemanticAnchor{ch1-1-3-1}
\noindent\textbf{1.3.1.}\quad
宇宙 $U$ を固定する。$U$ に属する $G$-次数付き圏を対象、次数付き関手を射と
する $U$-圏を $\operatorname{Grad}(G)$ と書く。同様に、$U$ に属する
$G^\circ$-圏を対象、$G^\circ$-関手を射とする $U$-圏を
$G^\circ\text{-}\operatorname{Cat}$ と書く。(1.2.1), (1.2.3), (1.2.5) から、
各 $G$-次数付き圏に正規化された切断系を選ぶことにより関手

\VerdierSemanticAnchor{ch1-1-3-1-eq1}
\begin{equation}
  \Psi\colon \operatorname{Grad}(G)
  \longrightarrow G^\circ\text{-}\operatorname{Cat}
  \tag{1.3.1.1}
\end{equation}
が定義されることが直ちに分かる。

[18] の結果は次のように解釈できる。

\VerdierSemanticAnchor{ch1-1-3-2}
\noindent\textbf{命題 1.3.2.}\quad
\emph{関手 $\Psi$ は圏同値である。}

\VerdierSemanticAnchor{ch1-1-3-2-quasi-inverse}
$\Psi$ の擬逆関手の記述は [18] にある。ここでは $G^\circ$-圏
$(\mathcal C,G^\circ,T,c)$ に対応する次数付き圏 $\mathcal C[G]$ を簡単に
記述する。まず $\mathcal C[G]$ の対象集合は $\mathcal C$ の対象集合である。
$X,Y$ を $\mathcal C$ の二対象とし、

\VerdierSemanticAnchor{ch1-1-3-2-eq1}
\begin{equation}
  \operatorname{Hom}^{g}(X,Y)
  =\operatorname{Hom}_{\mathcal C}(X,T(g)Y),
  \qquad g\in G,
  \tag{1.3.2.1}
\end{equation}

\VerdierSemanticAnchor{ch1-1-3-2-eq2}
\begin{equation}
  \operatorname{Hom}_{\mathcal C[G]}(X,Y)
  =\coprod_{g\in G}\operatorname{Hom}^{g}(X,Y)
  \tag{1.3.2.2}
\end{equation}
とおく。

\VerdierSemanticAnchor{ch1-1-3-2-composition}
次に $\mathcal C[G]$ における射の合成を記述する。$X,Y,Z$ を $\mathcal C$ の
三対象、$u\in\operatorname{Hom}^{g}(X,Y)$,
$v\in\operatorname{Hom}^{g'}(Y,Z)$（$g,g'\in G$）とする。関手 $T(g)$ は写像

\VerdierSemanticAnchor{ch1-1-3-2-eq3}
\begin{equation}
  \operatorname{Hom}^{g'}(Y,Z)
  \longrightarrow
  \operatorname{Hom}_{\mathcal C}
    \bigl(T(g)Y,T(g)T(g')Z\bigr)
  \tag{1.3.2.3}
\end{equation}
